% Part: propositional-logic
% Chapter: syntax-and-semantics
% Section: introduction

\documentclass[../../../include/open-logic-section]{subfiles}

\begin{document}

\olfileid{pl}{syn}{int}

\olsection{مقدمة}

في منطق القضايا، تُبنى !!{formula}s من
!!{propositional variable}s باستعمال الروابط القضوية
$\lnot$ و$\land$ و$\lor$ و$\lif$ و$\liff$. وحدسيًا، يمثّل
!!a{propositional variable}~$p$ جملةً أو قضيةً تكون صادقة أو كاذبة.
ومتى تحددت قيم الصدق لكل !!{propositional variable} وارد في
!!a{formula}، تحددت أيضًا قيم الصدق في أيٍّ من !!{formula}s التي تُبنى
من تلك المتغيرات باستعمال الروابط القضوية. ونقول إن منطق القضايا
\emph{دالّ صدقيًا}، لأن دلالته تُعطى بدوال تعتمد على قيم الصدق. وعلى
وجه الخصوص، نهمل في منطق القضايا أي تحديد إضافي للصدق والكذب، كأن يكون
شيء ما صادقًا
بالضرورة لا على نحو عارض، أو أن يكون معلوم الصدق، أو أن يكون صادقًا
الآن لا أن يكون قد صدق في الماضي أو سيصدق في المستقبل. ولا ننظر إلا
في قيمتي صدق: الصدق ($\True$) والكذب ($\False$)، فنستبعد من النقاش
إمكان ألا تكون العبارة صادقة ولا كاذبة، أو أن تكون نصف صادقة فحسب.
ونقتصر كذلك على الروابط التي تتحدد فيها قيمة صدق !!a{formula} مبنية
منها تحديدًا تامًا بقيم الصدق لأجزائها (لا بمعناها، مثلًا). وعلى وجه
الخصوص، ثمة خلاف في كون الجمل الشرطية في الإنجليزية دالّةً صدقيًا بهذا
المعنى. أما الشرط المادي~$\lif$ فدالّ صدقيًا؛ وتتناول
أنواع أخرى من المنطق جملًا شرطية ليست دالّةً صدقيًا.

ولكي نطوّر النظرية وما وراء النظرية لمنطق القضايا الدالّ صدقيًا، لا بد
أولًا من تعريف تركيب تعبيراته ودلالتها. وسنعرض طريقةً واحدةً لبناء
!!{formula}s من !!{propositional variable}s باستعمال الروابط. وثمة
تعريفات بديلة ممكنة. فستختار أنظمة أخرى رموزًا مختلفة، وتتخذ مجموعات
مختلفة من الروابط روابطَ أولية، وتستعمل الأقواس على نحو مختلف (أو لا
تستعملها أصلًا، كما في حالة ما يسمى بالتدوين البولندي). غير أن جميع
المقاربات تشترك في أن قواعد التكوين تعرّف مجموعة !!{formula}s
\emph{استقرائيًا}. وإذا صيغت هذه القواعد صياغة سليمة، لم يكن من الممكن،
في جوهر الأمر، أن يتولد كل تعبير بمقتضاها إلا على نحو واحد. وكون
التعريف الاستقرائي يفضي إلى تعبيرات \emph{ذات قراءة وحيدة} يعني أن
بوسعنا إسناد المعاني إلى هذه التعبيرات بالطريقة نفسها---أي بالتعريف
الاستقرائي.

إسناد المعنى إلى التعبيرات هو مجال الدلالة. والمفهوم المركزي في دلالة
منطق القضايا هو مفهوم الإرضاء في !!a{valuation}. يسند
!!^a{valuation}~$\pAssign{v}$ قيمتي الصدق $\True$ و$\False$ إلى
!!{propositional variable}s. ويعين أي !!{valuation} قيمة صدق
$\pValue{v}(!A)$ لكل !!{formula}~$!A$. وتكون !!^a{formula} مُرضاةً في
!!a{valuation}~$\pAssign{v}$ إذا وفقط إذا كانت
$\pValue{v}(!A) = \True$---ونكتب ذلك على الصورة $\pSat{v}{!A}$.
ويمكن أيضًا تعريف هذه العلاقة بالاستقراء على بنية~$!A$، وذلك باستعمال
دوال الصدق للروابط المنطقية لتعريف، مثلًا، متى تُرضى $!A \land !B$
بدلالة ما إذا كانت $!A$ مُرضاةً أم لا، وما إذا كانت~$!B$ مُرضاةً أم لا.

وعلى أساس علاقة الإرضاء $\pSat{v}{!A}$ للجمل يمكننا بعد ذلك تعريف
المفاهيم الدلالية الأساسية: التحصيل الحاصل، والاستلزام، وقابلية
الإرضاء. تكون !!^a{formula} تحصيلَ حاصل، $\Entails !A$، إذا أرضاها كل
!!{valuation}، أي إذا كانت $\pValue{v}(!A) = \True$ لأي~$\pAssign{v}$.
وتستلزمها مجموعة من !!{formula}s، $\Gamma \Entails !A$، إذا لزم، لكل
!!{valuation}، من إرضائه جميع !!{formula}s في~$\Gamma$ أن يرضي~$!A$
أيضًا. وتكون
مجموعة من !!{formula}s قابلةً للإرضاء إذا وجد !!{valuation} يرضي جميع
!!{formula}s فيها في آن واحد. ولأن !!{formula}s معرّفة استقرائيًا،
ولأن الإرضاء معرّف بدوره بالاستقراء على بنية !!{formula}s، يمكننا
استعمال الاستقراء لإثبات خصائص دلالتنا وربط المفاهيم الدلالية المعرّفة
بعضها ببعض.


\end{document}
